\subsection{FOURIER TRANSFORMS OF INFINITELY-DIFFERENTIABLE FUNCTIONS}

Recall (Chap. III, §3, no. 1, Def. 2) that U(G) denotes the algebra of distributions on G with support contained in \(\{e\}\) . The canonical injection of g into U(G) extends to an isomorphism from the enveloping algebra of the Lie algebra g to U(G) (loc. cit., no. 7, Prop. 25); from now on we identify these two algebras by this isomorphism. If f is an infinitely-differentiable complex function on G and if \(t \in \mathrm{U}(\mathrm{G})\) , we denote by \(L_{t}f\) and \(R_{t}f\) the functions on G defined by

\[
\mathrm{L} _ {t} f (g) = \langle \varepsilon_ {g} * t, f \rangle , \quad \mathrm{R} _ {t} f (g) = \langle t * \varepsilon_ {g}, f \rangle
\]

(cf.~loc. cit., no. 6). For all \(g \in G\) ,

\[
\mathrm{L} _ {t} \circ \gamma (g) = \gamma (g) \circ \mathrm{L} _ {t}, \quad \mathrm{R} _ {t} \circ \delta (g) = \delta (g) \circ \mathrm{R} _ {t}.
\]

Let \(u \in \hat{G}\) ; denote by \(E_u\) the space of \(u\) . The morphism of Lie groups \(u: G \to \mathbf{GL}(E_u)\) gives by differentiation a homomorphism of (real) Lie algebras \(\mathfrak{g} \to \text{End}(E_u)\) , hence a homomorphism of algebras, also denoted by \(u\) , from \(U(G)\) to \(\text{End}(E_u)\) . If \(t \in U(G)\) and if \(f\) is an infinitely-differentiable function on \(G\) , then

\[
u (\mathrm{L} _ {t} f) = u (f) u (t ^ {\vee}), \quad u (\mathrm{R} _ {t} f) = u (t ^ {\vee}) u (f),\tag{14}
\]

where \(t^{\vee}\) denotes the image of t under the principal anti-automorphism of U(G) (Chap. I, §2, no. 4); indeed, it suffices to verify this for \(t \in g\) , in which case it follows by differentiation from formulas (12) and (13) (cf.~Chap. III, §3, no. 7, Prop. 27).

For all \(u \in \hat{\mathbf{G}}\) , denote by \(\lambda(u)\) the highest weight of \(u\) (§7, no. 2, Th. 1), so \(u \mapsto \lambda(u)\) is a bijective map from \(\hat{\mathbf{G}}\) to the set \(\mathrm{X}_{++}\) of dominant elements of \(\mathrm{X}(\mathrm{T})\) .

Let \(\Gamma \in \mathrm{U}(\mathrm{G})\) be a Casimir element of \(\mathrm{G}\) (§7, no. 6); for all \(u \in \hat{\mathrm{G}}\) , the endomorphism \(u(\Gamma)\) of \(\mathrm{E}_u\) is a homothety, whose ratio we denote by \(\tilde{\Gamma}(u)\) , so we have a map \(u \mapsto \tilde{\Gamma}(u)\) from \(\hat{\mathrm{G}}\) to \(\mathbf{C}\) .

If \(\varphi\) and \(\psi\) are two functions on \(\hat{G}\) with positive real values, denote by `` \(\varphi \preccurlyeq \psi\) '' or `` \(\varphi(u) \preccurlyeq \psi(u)\) '' the relation ``there exists M \textgreater{} 0 such that \(\varphi(u) \leq M\psi(u)\) for all \(u \in \hat{G}\) ''; this is a pre-order relation on the set of functions on \(\hat{G}\) with positive real values.

PROPOSITION 2. Let \(m \mapsto \|m\|\) be a norm on the R-vector space \(\mathbf{R} \otimes \mathrm{X(T)}\) and \(\Gamma\) a Casimir element of G. Let \(\varphi\) be a function on \(\hat{\mathrm{G}}\) with positive real values.

\begin{enumerate}
\def\labelenumi{\alph{enumi})}
\tightlist
\item
  The following conditions are equivalent:
\end{enumerate}

\begin{enumerate}
\def\labelenumi{(\roman{enumi})}
\item
  There exists an integer \(n > 0\) such that \(\varphi(u) \preccurlyeq (\|\lambda(u)\| + 1)^n\) (resp. for every integer \(n > 0\) , we have \(\varphi(u) \preccurlyeq (\|\lambda(u)\| + 1)^{-n}\) ).
\item
  There exists an integer \(n > 0\) such that \(\varphi(u) \preccurlyeq (\tilde{\Gamma}(u) + 1)^n\) (resp. for every integer \(n > 0\) , we have \(\varphi(u) \preccurlyeq (\tilde{\Gamma}(u) + 1)^{-n}\) ).
\end{enumerate}

\begin{enumerate}
\def\labelenumi{\alph{enumi})}
\setcounter{enumi}{1}
\tightlist
\item
  If \(\mathrm{G}\) is semi-simple, conditions (i) and (ii) are also equivalent to:
\end{enumerate}

\begin{enumerate}
\def\labelenumi{(\roman{enumi})}
\setcounter{enumi}{2}
\tightlist
\item
  There exists an integer \(n > 0\) such that \(\varphi(u) \preccurlyeq d(u)^n\) (resp. for every integer \(n > 0\) , we have \(\varphi(u) \preccurlyeq d(u)^{-n}\) ).
\end{enumerate}

Note first of all that condition (i) is clearly independent of the choice of norm. Thus we can use the norm defined by the quadratic form \(Q_{\varGamma}\) associated to \(\varGamma\) (§7, no. 6, Prop. 4). Then

\[
0 \leq \tilde {\Gamma} (u) = \| \lambda (u) + \rho \| ^ {2} - \| \rho \| ^ {2},
\]

so \(\tilde{\Gamma}(u) + 1 \preccurlyeq (\|\lambda(u)\| + 1)^2 \preccurlyeq \tilde{\Gamma}(u) + 1\) , hence \(a\) .

Further, if G is semi-simple,

\[
\left\| \lambda (u) + \rho \right\| \preccurlyeq d (u) \preccurlyeq \left\| \lambda (u) + \rho \right\| ^ {\mathrm{N}}, \quad \text { where } \mathrm{N} = 1 / 2 (\dim \mathrm{G} - \dim \mathrm{T})
\]

(§7, no. 5, Cor. 1 of Th. 3), so \(\| \lambda(u) \| + 1 \preccurlyeq d(u) \preccurlyeq (\| \lambda(u) \| + 1)^{\mathrm{N}}\) , hence \(b\) ).

It follows from Prop. 2 that condition (i) is independent of the choice of maximal torus, chamber, and norm, and that condition (ii) is independent of the choice of Casimir element. A function \(\varphi\) satisfying conditions (i) and (ii) is said to be moderately increasing (resp. rapidly decreasing). The product of two moderately increasing functions is moderately increasing; the product of a moderately increasing function and a rapidly decreasing function is rapidly decreasing. If \(\varphi\) is rapidly decreasing, the family \((\varphi(u))_{u\in\hat{G}}\) is summable.

Examples. The function \(u \mapsto d(u)\) is moderately increasing (§7, no. 5, Cor. 1 of Th. 3); for any norm \(\| \cdot \|\) on \(\mathbf{R} \otimes \mathrm{X(T)}\) , the function \(u \mapsto \| \lambda(u) \|\) is moderately increasing. For any Casimir element \(\Gamma\) , the function \(u \mapsto \tilde{\Gamma}(u)\) is moderately increasing; more generally:

PROPOSITION 3. For all \(t \in \mathrm{U}(\mathrm{G})\) , the functions \(u \mapsto \| u(t) \|_{\infty}\) and \(u \mapsto \| u(t) \|_2\) on \(\hat{\mathrm{G}}\) are moderately increasing.

Since the product of two moderately increasing functions is moderately increasing, it suffices to prove this when \(t \in g\) : in that case the assertion follows from the Remark in §7, no. 6 and the inequality

\[
\| u (t) \| _ {2} \leq d (u) \| u (t) \| _ {\infty}.
\]

THEOREM 1. a) Let \(f\) be an infinitely-differentiable complex function on \(\mathbf{G}\) . Then the family \((f_u)_{u\in \hat{\mathbf{G}}}\) , where \(f_{u}(g) = \langle u(g)|u(f)\rangle\) , is uniformly summable on \(\mathbf{G}\) and, for all \(h\in \mathbf{G}\) ,

\[
f (h) = \sum_ {u \in \hat {\mathrm{G}}} \langle u (h) | u (f) \rangle = \sum_ {u \in \hat {\mathrm{G}}} d (u) \int_ {\mathrm{G}} f (g) \mathrm{Tr} (u (g h ^ {- 1}))   d g.
\]

\begin{enumerate}
\def\labelenumi{\alph{enumi})}
\setcounter{enumi}{1}
\tightlist
\item
  Let \(f\) be an integrable function on \(\mathbf{G}\) ; then \(f\) is equal almost everywhere to an infinitely-differentiable function if and only if the function \(u \mapsto \| u(f) \|_{\infty}\) is rapidly decreasing on \(\hat{\mathbf{G}}\) .
\end{enumerate}

Let \(f\) be an infinitely-differentiable function on \(G\) , and let \(\Gamma\) be a Casimir element for \(G\) ; by formula (14),

\[
\tilde {\Gamma} (u) ^ {n} u (f) = u (f) u (\Gamma) ^ {n} = u ((\mathrm{L} _ {\Gamma}) ^ {n} f)
\]

for all \(n \geq 0\) , and consequently

\[
\tilde {\Gamma} (u) ^ {n} \| u (f) \| _ {\infty} \leq \| (\mathrm{L} _ {\Gamma}) ^ {n} f \| _ {1} \leq \sup _ {g \in \mathrm{G}} | ((\mathrm{L} _ {\Gamma}) ^ {n} f) (g) |;\tag{15}
\]

thus, the function \(u \mapsto \|u(f)\|_{\infty}\) is indeed rapidly decreasing.

Conversely, let \(\mathrm{A} = (\mathrm{A}_{u})_{u \in \hat{\mathrm{G}}}\) be an element of \(\mathrm{F}(\hat{\mathrm{G}})\) such that the function \(u \mapsto \|A_{u}\|_{\infty}\) is rapidly decreasing. Put \(f_{u}(g) = \langle u(g)|A_{u} \rangle\) ; the function \(g \mapsto f_{u}(g)\) is analytic, hence infinitely-differentiable. By Chap. III, §3, no. 7, Prop. 27,

\[
(\mathrm{L} _ {x} f _ {u}) (g) = \langle u (g) u (x) | \mathrm{A} _ {u} \rangle
\]

for all \(x \in \mathfrak{g}\) . Let \(t \in \mathrm{U}(\mathrm{G})\) ; by the preceding formula,

\[
(\mathrm{L} _ {t} f _ {u}) (g) = \langle u (g) u (t) | \mathrm{A} _ {u} \rangle
\]

and consequently

\[
\begin{array}{r} | (\mathrm{L} _ {t} f _ {u}) (g) | = | \langle u (g) u (t) | \mathrm{A} _ {u} \rangle \leq d (u) ^ {2} \| u (t) \| _ {\infty} \| u (g) \| _ {\infty} \| \mathrm{A} _ {u} \| _ {\infty} \\ = d (u) ^ {2} \| u (t) \| _ {\infty} \| \mathrm{A} _ {u} \| _ {\infty}. \end{array}
\]

Since \(d(u)\) and \(\|u(t)\|_{\infty}\) are moderately increasing (Prop. 3) and \(\|A_{u}\|_{\infty}\) is rapidly decreasing, the function \(u \mapsto \sup_{g} |(\mathrm{L}_{t}f_{u})(g)|\) is rapidly decreasing; thus, the family \((\mathrm{L}_{t}f_{u})_{u \in \hat{\mathrm{G}}}\) is uniformly summable. It follows \(^{7}\) that the sum of the family \((f_{u})\) is an infinitely-differentiable function on G, whose Fourier cotransform is \((\mathrm{A}_{u})\) , hence the theorem.

Denote by \(\mathcal{S}(\hat{\mathrm{G}})\) the vector subspace of \(\mathrm{L}^2 (\hat{\mathrm{G}})\) consisting of the families \(\mathrm{A} = (\mathrm{A}_u)_{u\in \hat{\mathrm{G}}}\) such that the function \(u\mapsto \| \mathrm{A}_u\|_\infty\) is rapidly decreasing on \(\hat{\mathrm{G}}\) . It follows from the theorem that the maps \(\overline{\mathcal{F}}:f\mapsto (u(f))_{u\in \hat{\mathrm{G}}}\) and \(\mathcal{F}\colon \mathrm{A}\mapsto \sum_{u\in \hat{\mathrm{G}}}\langle u(g)|\mathrm{A}_u\rangle\) induce inverse isomorphisms between the complex vector spaces \(\mathcal{C}^\infty (\mathrm{G};\mathbf{C})\) and \(\mathcal{S}(\hat{\mathrm{G}})\) . Give the space \(\mathcal{C}^\infty (\mathrm{G};\mathbf{C})\) the topology of uniform \(\mathrm{C}^\infty\) -convergence (§6, no. 4) which can be defined by the family of semi-norms \(f\mapsto \sup_{g\in \mathrm{G}}|\mathrm{L}_t f(g)|\) for \(t\in \mathrm{U}(\mathrm{G})\) , and the space \(\mathcal{S}(\hat{\mathrm{G}})\) the topology defined by the sequence of semi-norms \(p_n:\mathrm{A}\mapsto \sup_{u\in \hat{\mathrm{G}}}(\tilde{\Gamma} (u) + 1)^n\| \mathrm{A}_u\|_\infty\) . Formula (15) of the preceding proof shows that \(\overline{\mathcal{F}}\) is continuous. Let \(t\in \mathrm{U}(\mathrm{G})\) , and let \(\mathrm{A} = (\mathrm{A}_u)_{u\in \hat{\mathrm{G}}}\) be an element of \(\mathcal{S}(\hat{\mathrm{G}})\) ; put \(f_{n}(g) = \langle u(g)|\mathrm{A}_{u}\rangle\) . Let \(p\) be an integer such that \(\sum_{u\in \hat{\mathrm{G}}} \tilde{\Gamma} (u)^{-p} = \mathrm{M} < \infty\) . By the preceding proof, there exists a positive integer \(m\) such that, for all \(g\in \mathrm{G}\) ,

\[
| (\mathrm{L} _ {t} f _ {u}) (g) | \leq d (u) ^ {2} \| u (t) \| _ {\infty} \| \mathrm{A} _ {u} \| _ {\infty} \leq m. (1 + \tilde {\Gamma} (u)) ^ {m} \tilde {\Gamma} (u) ^ {- p} \| \mathrm{A} _ {u} \| _ {\infty}
\]

so \(|(\mathrm{L}_t\mathcal{F}(\mathrm{A}))(g)| \leq m\mathrm{M}p_m(\mathrm{A})\) ; this proves that \(\mathcal{F}\) is continuous. Consequently:

COROLLARY. The maps \(\overline{\mathcal{F}}: f \mapsto (u(f))_{u \in \hat{\mathrm{G}}} \text{ and } \mathcal{F}: \mathrm{A} \mapsto \sum_{u \in \hat{\mathrm{G}}} \langle u(g) | \mathrm{A}_u \rangle\) induce inverse isomorphisms between the topological vector spaces \(\mathcal{C}^\infty(\mathrm{G}; \mathbf{C})\) and \(\mathcal{S}(\hat{\mathrm{G}})\) .

